\documentclass[14pt]{extarticle}
\usepackage{fontspec}
\usepackage{polyglossia}
\usepackage{amsmath}
\usepackage{amssymb}
\usepackage{xcolor}
\usepackage{fancyhdr}
\usepackage{graphicx}
\usepackage{listings}
\usepackage[most]{tcolorbox}
% \usepackage{setspace}
\usepackage{pifont}
\usepackage{enumitem}
\usepackage{titlesec}
\usepackage[bottom]{footmisc}
\usepackage{titling}
\usepackage{minted}
\usepackage{etoolbox}
\usepackage{array}
\usepackage{longtable} % page-breakable tables (must load before polyglossia/bidi RTL setup)
\usepackage{extsizes}
\usepackage{float}
\usepackage{booktabs}
\usepackage{hyperref}
\usepackage[onehalfspacing]{setspace}
\newif\ifwhiteblack
\ifdefined\setwhiteblack
  \setwhiteblack
\fi

\ifwhiteblack
  \usepackage[a4paper, landscape, top=1.25cm, bottom=1.25cm, left=1.25cm, right=1.25cm, headheight=1.5cm, headsep=0.5cm]{geometry}
\else
  \usepackage[a4paper, top=2.7cm, bottom=2.7cm, left=2.5cm, right=2.5cm, headheight=1.5cm, headsep=0.5cm]{geometry}
\fi
\usepackage{tikz}
\usepackage{datetime2}

\usetikzlibrary{automata, positioning, arrows.meta, arrows}
\providecommand{\LTcaptype}{table}

\newfontfamily\emoji{Segoe UI Emoji}

\pagestyle{fancy}

\setmainlanguage[numerals=western]{arabic}
\setotherlanguage{english}
\newfontfamily\arabicfont[Script=Arabic]{Calibri}
\newfontfamily\arabicfonttt[Script=Arabic]{Courier New}



\ifwhiteblack
\lstset{
  language=[Sharp]C,
  numbers=left,
  stepnumber=1,
  numbersep=8pt,
  frame=single,
  basicstyle=\ttfamily\small,
  keywordstyle=\bfseries,
  stringstyle=\ttfamily,
  commentstyle=\itshape
}
\else
\lstset{
  language=[Sharp]C,
  numbers=left,
  stepnumber=1,
  numbersep=8pt,
  frame=single,
  basicstyle=\ttfamily\small,
  keywordstyle=\color{blue},
  stringstyle=\color{red},
  commentstyle=\color{green!50!black}
}
\fi

\newif\ifdetailed
\ifdefined\setdetailed
  \setdetailed
\fi

\newif\ifwithsols
\ifdefined\setwithsols
  \setwithsols
\fi

\ifwhiteblack
  \AtBeginDocument{
    \LARGE
    \usemintedstyle{bw}
  }
\fi

% Default minted look: compact code that wraps long lines.
% Not applied in black-and-white landscape mode, which already uses \LARGE.
\ifwhiteblack\else
  \setminted{fontsize=\small, breaklines=true}
\fi

% unified tcolorboxes for articles
\ifwhiteblack
\tcbset{colback=white, colframe=black, fonttitle=\bfseries, boxrule=1pt}
\newtcolorbox{boxDef}[1][]{title={{\emoji📘} تعريف\ifx\\#1\\\else ~#1\fi :}}
\newtcolorbox{boxExercise}[1][]{title={{\emoji🧩} تمرين\ifx\\#1\\\else ~#1\fi :}}
\newtcolorbox{boxExample}[1][]{title={{\emoji📝} مثال\ifx\\#1\\\else ~#1\fi :}}
\newtcolorbox{boxNote}[1][]{title={{\emoji✨} ملاحظة\ifx\\#1\\\else ~#1\fi :}}
\newtcolorbox{boxAttention}[1][]{title={{\emoji🔔} تنبيه\ifx\\#1\\\else ~#1\fi :}}
\newtcolorbox{boxWarning}[1][]{title={{\emoji⚡} ملاحظة هامة\ifx\\#1\\\else ~#1\fi :}}
\newtcolorbox{boxSolution}[1][]{title={{\emoji✅} حل\ifx\\#1\\\else ~#1\fi :}}
\newtcolorbox{boxSymbol}[1][]{title={{\emoji🔣} رمز\ifx\\#1\\\else ~#1\fi :}}
\newtcolorbox{boxHint}[1][]{title={{\emoji💡} تلميح\ifx\\#1\\\else ~#1\fi :}}
\else
\tcbset{colback=white, colframe=black, fonttitle=\bfseries, boxrule=0.8pt}
\newtcolorbox{boxDef}[1][]{colback=blue!5!white,colframe=blue!75!black,
  title={{\emoji📘} تعريف\ifx\\#1\\\else ~#1\fi :}}
\newtcolorbox{boxExercise}[1][]{colback=cyan!5!white,colframe=cyan!70!black,
  title={{\emoji🧩} تمرين\ifx\\#1\\\else ~#1\fi :}}
\newtcolorbox{boxExample}[1][]{colback=yellow!5!white,colframe=orange!90!black,
  title={{\emoji📝} مثال\ifx\\#1\\\else ~#1\fi :}}
\newtcolorbox{boxNote}[1][]{colback=gray!10!white,colframe=black,
  title={{\emoji✨} ملاحظة\ifx\\#1\\\else ~#1\fi :}}
\newtcolorbox{boxAttention}[1][]{colback=magenta!10!white,colframe=magenta!80!black,
  title={{\emoji🔔} تنبيه\ifx\\#1\\\else ~#1\fi :}}
\newtcolorbox{boxWarning}[1][]{colback=red!5!white,colframe=red!75!black,
  title={{\emoji⚡} ملاحظة هامة\ifx\\#1\\\else ~#1\fi :}}
\newtcolorbox{boxSolution}[1][]{colback=green!5!white,colframe=green!60!black,
  title={{\emoji✅} حل\ifx\\#1\\\else ~#1\fi :}}
\newtcolorbox{boxSymbol}[1][]{colback=purple!5!white,colframe=purple!70!black,
  title={{\emoji🔣} رمز\ifx\\#1\\\else ~#1\fi :}}
\newtcolorbox{boxHint}[1][]{colback=teal!5!white,colframe=teal!60!black,
  title={{\emoji💡} تلميح\ifx\\#1\\\else ~#1\fi :}}
\fi


\ifwhiteblack
\tcbset{simplecode/.style={ colback=white, colframe=black, boxrule=1pt, arc=2pt, left=4pt,right=4pt,top=4pt,bottom=4pt}}
\else
\tcbset{simplecode/.style={ colback=gray!5, colframe=black!50, boxrule=0.4pt, arc=2pt, left=4pt,right=4pt,top=4pt,bottom=4pt}}
\fi
% Compact code boxes: tighter line spacing (the preamble uses onehalfspacing) and small vertical gaps.
% Must come AFTER the simplecode definition above.
\tcbset{simplecode/.append style={fontupper=\setstretch{0.95}, top=4pt, bottom=4pt, before skip=6pt, after skip=6pt}}
\newenvironment{boxCode}{\begin{tcolorbox}[simplecode]}{\end{tcolorbox}}

% Code cell used inside tables (keeps code left-to-right inside RTL tables).
% Use inside a fixed-width p{...} column: \begin{pycell}\begin{minted}...\end{minted}\end{pycell}
\newtcolorbox{pycell}{frame empty,colback=white,boxsep=0pt,left=2pt,right=2pt,top=2pt,bottom=2pt,width=\linewidth}

\newcolumntype{C}[1]{>{\centering\arraybackslash}p{#1}}

\makeatletter
\DTMnewdatestyle{dotted}{%
  \renewcommand{\DTMdisplaydate}[4]{\two@digits{##3}.\two@digits{##2}.##1}%
  \renewcommand{\DTMDisplaydate}{\DTMdisplaydate}%
}
\makeatother
\DTMsetdatestyle{dotted}

% redefine spaces after titles
\makeatletter
\renewcommand{\@maketitle}{%
  \begin{center}
    {\huge \bfseries \@title \par}%
    \vskip 0.2em % space between title and author
    {\large \@author \par}%
    % \vskip 0.2em % space between author and date
    % {\normalsize \@date \par}%
    {\small \hfill \begin{english}\DTMtoday\end{english} \par}%
  \end{center}
}
\makeatother

\fancyhf{} % clear default
\fancypagestyle{plain}{
  \fancyhf{}
  \fancyhead[L]{مدرسة التسامح الشاملة}
  % \fancyhead[L]{\includegraphics[height=1cm]{../../../images/logoTasamoh.png}}
  \fancyhead[R]{الأستاذ محمود اغبارية}
  \fancyfoot[C]{\thepage}
}

\fancyhead[L]{مدرسة التسامح الشاملة}
\fancyhead[R]{الأستاذ محمود اغبارية}
\fancyfoot[C]{\thepage}

\setcounter{tocdepth}{3} % only section subsection and subsubsection in TOC

\newcommand{\importpdfpage}[2]{%
    \thispagestyle{empty}
    \newgeometry{left=0cm,right=0cm,top=0cm,bottom=0cm}
    \includegraphics[page=#2,width=0.95\paperwidth,height=0.95\paperheight,keepaspectratio]{#1}
    \restoregeometry
}

\newcommand{\insertFullImg}[1]{%
    \noindent
    \makebox[\textwidth][c]{
        \includegraphics[width=0.9\paperwidth,keepaspectratio]{#1}%
    }%
}%



\title{أكاديمية فرسان التنانين — المرحلة الثالثة: سجلّ الوادي وتعدد الأشكال}

\begin{document}

\maketitle
\thispagestyle{fancy}

\begin{boxWarning}
\textbf{تعليمات هامة للعمل:}
تعتمد هذه المرحلة على الفئات التي بنيتموها في المراحل السابقة (\textenglish{BasicDragon} و \textenglish{FireDragon} و \textenglish{HibernatingDragon} و \textenglish{AlphaDragon}) — حدّثوا الفئات الموجودة ولا تبدأوا من الصفر. \\
في \textbf{كل تمرين}، يجب تحديث مقطع الفحص داخل العملية الرئيسية \textenglish{Main} وتشغيله للتأكد من صحة عمل الفئات والعمليات الجديدة التي كتبتموها.
\end{boxWarning}

% ==============================================================
\section*{المرحلة الثالثة: سجلّ الوادي وتعدد الأشكال}
% ==============================================================

\begin{enumerate}[label=\textbf{تمرين \arabic*:}, wide=0pt, leftmargin=*, itemsep=2.2em]

    \item \textbf{فئة سجلّ الوادي (\textenglish{DragonRegistry})} \\
    ابنوا فئة خدمات باسم \textenglish{DragonRegistry} تدير كل تنانين الوادي:

    {\centering
    \renewcommand{\arraystretch}{1.2}
    \small
    \begin{tabular}{|>{\raggedleft\arraybackslash}p{0.4\textwidth}|>{\raggedright\arraybackslash}p{0.5\textwidth}|}
    \hline
    \multicolumn{2}{|c|}{\textbf{الفئة: \textenglish{DragonRegistry}}} \\
    \hline
    صفات: مصفوفة تحوي تنانين من جميع الأنماط (\textenglish{FireDragon}، \textenglish{HibernatingDragon}، \textenglish{AlphaDragon})، وعدّاد لعدد التنانين الفعليين. & \textenglish{private BasicDragon[] dragons;} \newline \textenglish{private int count;} \\
    \hline
    عملية بنائية تتلقى حجم المصفوفة وتنشئها. & \textenglish{public DragonRegistry(int maxDragons)} \\
    \hline
    عملية مساعدة تتلقى رقم تنين وتعيد مؤشره (مكانه) في المصفوفة أو $-1$. & \textenglish{private int FindIndex(int dragonNum)} \\
    \hline
    عملية تضيف التنين في أول مكان فارغ، وتعيد \textenglish{true} إذا أُضيف و \textenglish{false} إذا كانت المصفوفة ممتلئة أو كان المُعطى \textenglish{null}. & \textenglish{Add(BasicDragon d)} \\
    \hline
    عملية تعيد التنين صاحب الرقم (أو \textenglish{null}) بالاستعانة بـ \textenglish{FindIndex}. & \textenglish{Find(int dragonNum)} \\
    \hline
    عملية تحذف التنين وتُزيح الحدود التالية إلى اليسار (بلا ثغرات في الوسط) وتعيد \textenglish{bool}. & \textenglish{Remove(int dragonNum)} \\
    \hline
    عملية تعيد عدد التنانين الموجودين فعلياً، وعملية تعيد نصاً يصف كل التنانين في السجل. & \textenglish{Count(), ToString()} \\
    \hline
    \end{tabular}
    \par} %

    \vspace{2em}
    \noindent \textbf{الفحص في \textenglish{Main}:} \\
    أنشئوا سجلاً بحجم $6$، وأضيفوا إليه $5$ تنانين من الأنماط الثلاثة، وجرّبوا \textenglish{Find} و \textenglish{Remove} و \textenglish{Count} وطباعة السجل.

    \item \textbf{العمليات القابلة لإعادة التعريف (\textenglish{virtual} و \textenglish{override}): \textenglish{IsInDanger} و \textenglish{DangerLevel}}
    \begin{enumerate}[label=\textbf{(\arabic*)}]
        \item أضيفوا إلى المشروع عملية عامة قابلة لإعادة التعريف تعريفها \textenglish{public virtual bool IsInDanger()}:
        \begin{itemize}
            \item في \textenglish{BasicDragon}: تعيد \textenglish{false} دائماً.
            \item في \textenglish{FireDragon}: تعيد \textenglish{true} إذا كان التنين يستخدم جمر الطوارئ ($\text{energy} < 0$).
            \item في \textenglish{HibernatingDragon}: تعيد \textenglish{true} إذا كان التنين بلا طاقة ($\text{energy} \le 0$).
            \item في \textenglish{AlphaDragon}: تعيد \textenglish{true} إذا كان $\text{emberReserve} > 0$ واستخدم القائد $90\%$ أو أكثر من جمره ($\text{energy} \le -0.9 \cdot \text{emberReserve}$).
        \end{itemize}
        \item أضيفوا إلى \textenglish{DragonRegistry} عملية \textenglish{DangerDragons()} تعيد مصفوفة (\textenglish{BasicDragon[]}) بالتنانين التي في خطر.
        \item أضيفوا عملية قابلة لإعادة التعريف \textenglish{public virtual string DangerLevel()} تعيد \textenglish{"Safe"} أو \textenglish{"Warning"} أو \textenglish{"Critical"}:
        \begin{itemize}
            \item في \textenglish{BasicDragon}: دائماً \textenglish{"Safe"}.
            \item في \textenglish{FireDragon}: إذا لم تكن الطاقة سالبة $\leftarrow$ \textenglish{"Safe"}. وإذا كانت سالبة: إن استُهلك نصف جمر الطوارئ أو أكثر ($\text{energy} \le -0.5 \cdot \text{emberReserve}$) $\leftarrow$ \textenglish{"Critical"}، وإلا $\leftarrow$ \textenglish{"Warning"}.
            \item في \textenglish{HibernatingDragon}: \textenglish{"Critical"} إذا كانت الطاقة $0$ أو أقل، وإلا \textenglish{"Safe"}.
            \item في \textenglish{AlphaDragon}: إذا كان القائد في خطر (\textenglish{IsInDanger()}) $\leftarrow$ \textenglish{"Critical"}. وإلا إذا كانت طاقته سالبة $\leftarrow$ \textenglish{"Warning"}. وإلا افحصوا أعضاء القطيع: إذا كان أحدهم \textenglish{"Critical"} $\leftarrow$ \textenglish{"Warning"}، وإلا $\leftarrow$ \textenglish{"Safe"}.
        \end{itemize}
    \end{enumerate}
    \textbf{الفحص في \textenglish{Main}:} \\
    اطبعوا \textenglish{DangerLevel()} لكل تنين في السجل، ثم استدعوا \textenglish{DangerDragons()} واطبعوا النتيجة.

    \item \textbf{أذونات الوصول \textenglish{protected} والكلمة \textenglish{base} وسلسلة \textenglish{ToString}}
    \begin{itemize}
        \item أضيفوا إلى \textenglish{BasicDragon} صفة محمية \textenglish{protected string dragonType} تُحدَّد قيمتها في كل فئة وارثة (\textenglish{"Fire"} أو \textenglish{"Hibernating"} أو \textenglish{"Alpha"}).
        \item غيّروا \textenglish{ToString()} في \textenglish{BasicDragon} لتعيد:
        \begin{center}
        \textenglish{\$"Type:\{dragonType\}, Num:\{dragonNumber\}, Rider:\{riderID\}, Energy:\{energy\}"}
        \end{center}
        \item في كل فئة وارثة، أعيدوا تعريف \textenglish{ToString()} بحيث تستدعي \textenglish{base.ToString()} وتضيف الصفات الخاصة بها.
    \end{itemize}
    \textbf{الفحص في \textenglish{Main}:} \\
    اطبعوا السجل وتأكدوا من صيغة \textenglish{ToString()} الجديدة لجميع التنانين.

\clearpage

    \item \textbf{عمليات خاصة بفئة فرعية واحدة فقط}
    \begin{itemize}
        \item أضيفوا إلى \textenglish{HibernatingDragon} عملية \textenglish{IsEmpty()} تعيد \textenglish{true} إذا كان التنين بلا طاقة ($\text{energy} \le 0$)، واستخدموها داخل \textenglish{IsInDanger()}.
        \item أضيفوا إلى \textenglish{FireDragon} عملية \textenglish{RescueSuggestion()} تعيد سلسلة نصية: إذا كانت طاقة التنين سالبة تقترح طلب تحويل طاقة من قائد قطيع بمقدار الدَّين مقرّباً لأعلى (\textenglish{Math.Ceiling(-energy)})، وإلا تعيد \textenglish{"No help needed"}.
    \end{itemize}
    \textbf{الفحص في \textenglish{Main}:} \\
    استدعوا العمليتين على كائنات من نمطيهما في \textenglish{Main}.

    \item \textbf{العامل \textenglish{is} والتحويل لأسفل (\textenglish{PrintSuggestions})} \\
    أضيفوا إلى \textenglish{DragonRegistry} عملية \textenglish{PrintSuggestions()} تمسح كل تنانين السجل وتفعل ما يلي:
    \begin{itemize}
        \item إذا كان التنين \textenglish{HibernatingDragon} وكان فارغاً (\textenglish{IsEmpty()}) $\leftarrow$ تطبع تحذيراً برقم التنين.
        \item إذا كان \textenglish{FireDragon} وكان في خطر (\textenglish{IsInDanger()}) $\leftarrow$ تطبع \textenglish{RescueSuggestion()}.
        \item إذا كان \textenglish{AlphaDragon}: أضيفوا له عملية \textenglish{RescueSuggestion()}، واطبعوها فقط إذا كان القائد في خطر \textbf{أو} كان أحد أعضاء قطيعه في مستوى \textenglish{"Critical"}.
    \end{itemize}
    \textbf{الفحص في \textenglish{Main}:} \\
    استدعوا \textenglish{PrintSuggestions()} على السجل وتأكدوا من طباعة الرسائل المتوقعة.

    \item \textbf{العامل \textenglish{as} ومقارنة الأساليب (\textenglish{PrintSuggestions2})} \\
    أعيدوا كتابة \textenglish{PrintSuggestions} في نسخة جديدة باسم \textenglish{PrintSuggestions2()} باستخدام العامل \textenglish{as} بدل \textenglish{is}، وقارنوا في تعليق بين \textenglish{is} و \textenglish{as} والتحويل المباشر. \\
    \textbf{الفحص في \textenglish{Main}:} \\
    شغّلوا \textenglish{PrintSuggestions()} و \textenglish{PrintSuggestions2()} وتأكدوا من تطابق المخرجات.

    \item \textbf{التغذية الجماعية (\textenglish{FeedAll} و \textenglish{TotalEnergy})} \\
    أضيفوا إلى \textenglish{DragonRegistry}:
    \begin{itemize}
        \item العملية \textenglish{FeedAll(double amount)}: تتلقى مقدار طاقة وتُعبّئ به كل تنانين السجل، وتعيد عدد التنانين التي نجحت تعبئتها.
        \item العملية \textenglish{TotalEnergy()}: تعيد مجموع طاقة كل التنانين في السجل.
    \end{itemize}
    \textbf{الفحص في \textenglish{Main}:} \\
    اطبعوا \textenglish{TotalEnergy()} قبل وبعد استدعاء \textenglish{FeedAll(25)} و \textenglish{FeedAll(-10)}.

    \item \textbf{الوراثة من \textenglish{Object} — العملية \textenglish{Equals}}
    \begin{itemize}
        \item أعيدوا تعريف العملية \textenglish{Equals(object obj)} في \textenglish{BasicDragon} بحيث يعتبر التنينان متساويين إذا تطابقت لديهما قيم: رقم المملكة ورقم الكهف ورقم التنين.
        \item عدّلوا \textenglish{Add} في \textenglish{DragonRegistry} بحيث لا تضيف تنيناً موجوداً مسبقاً في السجل (بحسب \textenglish{Equals}).
    \end{itemize}
    \textbf{الفحص في \textenglish{Main}:} \\
    اختبروا \textenglish{Equals} وحاولوا إضافة تنين مكرر إلى السجل.

    \item \textbf{دمج سجلّين (\textenglish{Merge})} \\
    أضيفوا إلى \textenglish{DragonRegistry} عملية \textenglish{Merge(DragonRegistry other)} تتلقى سجلاً آخر وتعيد \textbf{سجلاً جديداً} يحوي تنانين السجلّين معاً \textbf{دون تكرار} (باستخدام \textenglish{Add}). \\
    \textbf{الفحص في \textenglish{Main}:} \\
    أنشئوا سجلّين بينهما تنانين مشتركة، ادمجوهما عبر \textenglish{Merge}، واطبعوا السجل المدموج.

\end{enumerate}

\vspace{1cm}
\begin{flushleft}
أرجو لكم وقتًا ممتعًا.

الأستاذ محمود اغبارية.
\end{flushleft}

\end{document}
